\documentclass[12pt]{article}%
\usepackage{graphicx, verbatim}
\usepackage{amsmath}
\usepackage{amsfonts}
\usepackage{amssymb}
\usepackage{latexsym}
\usepackage{enumerate}
\usepackage[spanish, provide=*]{babel}
\usepackage{palatino}
\usepackage[utf8]{inputenc}


\setlength{\headsep}{-2cm} \setlength{\oddsidemargin}{-0.9cm}
\setlength{\evensidemargin}{1.5cm} \setlength{\textheight}{25cm}
\setlength{\textwidth}{17.7cm} \setlength{\parindent}{0cm}
\parindent=0mm

\pagestyle{empty}

\def\R{\mathbb{R}}
\def\Z{\mathbb{Z}}
\def\C{\mathbb{C}}
\def\N{\mathbb{N}}
\def\A{\mathcal{Z}}
\def\V{\mathcal V}
\def\W{\mathcal W}
\def\F{\mathcal F}
\def\a{\alpha}
\def\vf{\varphi}
\def\re{\operatorname{Re}}
\def\log{\operatorname{Log}}
\def\adj{\operatorname{adj}}
\def\img{\operatorname{Im}}
\def\nu{\operatorname{Nu}}
\newtheorem{exercise}{Ejercicio}
\newcommand{\bej}{\begin{exercise}\rm}
\newcommand{\fej}{\end{exercise}}

% \renewcommand{\labelenumi}{\bf{(\numb{enumi})}}
\renewcommand{\labelenumii}{(\alph{enumii})}

\begin{document}

\pagestyle{empty}%

\label{t1}


\bigskip

%\hfill\textsc{\textbf{Tema }}

\textsc{Nombre y apellido:}

\vskip0.2cm

\textsc{DNI y Libreta Universitaria:}

\vskip0.2cm

\textsc{Carrera:}

\vskip0.4cm


\def\espac{\LARGE$\phantom{\displaystyle\sum\, \,}$}

\begin{tabular}[c]{|c | c | c | c| c|}
\hline
1 & 2 & 3 & 4 & Calif. \\
\hline
\espac& \espac & \espac & \espac & \espac \\
\hline
\end{tabular}

\bigskip
\medskip

\noindent

\hrulefill\ \vskip8pt

\begin{large}\centerline{{\textbf{Cálculo Numérico / Elementos de Cálculo Numérico}}}\end{large}

\vskip.2cm

\begin{large}\centerline{{\textbf{Primer Parcial - 6 de Mayo de 2026}}}\end{large}

\hrulefill

\vskip 0.5cm

\begin{exercise}[Cuadrados mínimos ponderados]
Se tienen mediciones con diferente grado de confiabilidad, y se desea ajustar un modelo lineal a los datos. Para esto, se asignan pesos $w_i>0$ a cada medición, y se plantea el problema de cuadrados mínimos ponderados
\[\min_{x\in\R^n} \sum_{i=1}^m w_i (a_i^Tx - b_i)^2,\]
donde $a_i^T$ es la $i$-ésima fila de $A\in\R^{m\times n}$ y $b\in\R^m$ es el vector de datos.

\begin{enumerate}[(a)]
\item Obtenga las ecuaciones normales para este problema. (Sugerencia: verifique que el problema se puede escribir como un problema de cuadrados mínimos estándar $\min_{x\in\R^n} \|\tilde{A}x - \tilde{b}\|_2^2$ con $\tilde{A}\in\R^{m\times n}$ y $\tilde{b}\in\R^m$ definidos en función de $A$, $b$ y $w_i$.)  
\item Proponga un método numérico para resolver el problema de cuadrados mínimos ponderados basado en la descomposición QR. 
\end{enumerate}
\end{exercise}

\begin{exercise}[Integración numérica con peso singular]
Se desean calcular integrales de la forma 
\[ \int_0^1 \ln(x) f(x) dx, \]
donde $f\in C^2([0,1])$ es una función dada. 

\begin{enumerate}[(a)]
    \item ¿Qué ocurre si se intenta aplicar la regla del trapecio compuesta con nodos equiespaciados en $x$? Aplique el cambio de variables $x=t^n$ con $n\in\mathbb{N}$. ¿A partir de qué valor de $n$ el integrando es una función $C^2$? 
    \item Para el valor de $n$ hallado, escribir la aproximación de la nueva integral utilizando la regla de trapecios compuesta con nodos equiespaciados en $t$ ¿Cuál es el orden de convergencia de este método?
\end{enumerate}

\end{exercise}

\begin{exercise}[Interpolación trigonométrica y FFT]
Dados $n$ valores $f_0,\ldots,f_{n-1}\in\C$, se desea encontrar coeficientes $c_0,\ldots,c_{n-1}\in\C$ tales que el polinomio trigonométrico
\[
p(t) = \sum_{j=0}^{n-1} c_j\,(e^{2\pi i t})^j
\]
interpole los datos en cierto conjunto de nodos $t_k \in [a,b]$, es decir, $p(t_k) = f_k$ para $k=0,\ldots,n-1$.

\begin{enumerate}[(a)]
\item Fijados $\{ t_k \}$ cualesquiera, escriba el sistema lineal que hay que resolver para obtener los coeficientes $\{c_j \}$ a partir de $\{ f_k \}$.
\item Considere ahora el polinomio trigonométrico reescalado en $[a,b]$
\[
g(t)=\sum_{j=0}^{n-1} d_j\,e^{2\pi i j\frac{t-a}{b-a}}.
\]
Si los nodos son equiespaciados $t_k=a+\frac{b-a}{n}k$, $k=0,\ldots,n-1$, y los datos cumplen $g(t_k)=f_k$, escriba el sistema lineal para los coeficientes $\{d_j\}$ y explique cómo resolverlo usando FFT/IFFT.
\end{enumerate}
\end{exercise}

\begin{exercise}[Transformada rápida de Hadamard]
Se define la familia de matrices de Hadamard $H_m\in\R^{2^m\times 2^m}$ de forma recursiva por
\[
H_0 = (1), \qquad
H_m = \begin{pmatrix} H_{m-1} & H_{m-1} \\ H_{m-1} & -H_{m-1} \end{pmatrix}, \quad m\ge 1.
\]
Dado $z\in\R^{n}$ con $n=2^m$, se llama \emph{transformada de Hadamard} de $z$ al vector $\hat{z} = H_m z$.

\begin{enumerate}[(a)]
\item Escribir el pseudocódigo de un método para calcular $\hat{z} = H_m z$ con costo computacional de $O(n \mathrm{log}(n))$.

\item Muestre que $H_m^T = H_m$ y que $H_m^2 = nI$, donde $n = 2^m$ (Sugerencia: puede utilizar inducción y productos de matrices a bloques, o al menos verificar explícitamente para $m=1,2$). 

\item Usando las propiedades de la matriz $H_m$ obtenidas en el item anterior, obtenga que $\mathrm{cond}_2(H_m)=1$. ¿Qué implica este resultado sobre la estabilidad numérica del cálculo de $\hat{z}$ ante perturbaciones en $z$?
\end{enumerate}
\end{exercise}


\vskip 2.5cm
\begin{center}
    \textit{En todos los casos debe justificar sus respuestas.}
\end{center}


\end{document}
